\subsubsection*{Data Sources}

The modeling pipeline integrates satellite-derived variables, weather data, terrain information, and in-situ soil moisture observations. The target variable throughout the study is soil moisture measured at a depth of 5 cm.

\subsubsection*{Satellite Data}
Satellite-based predictors were collected using Google Earth Engine \cite{gee}. The data sources include Sentinel-1, Sentinel-2, MODIS, SMAP, and a digital elevation model (DEM).

Sentinel-1 synthetic aperture radar (SAR) data \cite{sentinel1} were used to extract VV and VH backscatter in both decibel and linear representations. 
Sentinel-2 Level-2A surface reflectance data \cite{sentinel2} provided spectral bands \(B2\), \(B3\), \(B4\), \(B8\), \(B11\), and \(B12\). MODIS products \cite{modis} were used to obtain land surface temperature (LST) and NDVI. Terrain-related variables were derived from the SRTM DEM \cite{srtm}, including elevation, slope, and aspect. 
SMAP data \cite{smap} provided AM and PM soil moisture estimates along with retrieval quality flags. SoilGrids 2.0 \cite{Poggioetal2021} was used to collect the soil type. 

To improve robustness and reduce noise, satellite values were aggregated over a spatial buffer around each station location. A buffer radius of 1000 m was used. The effective spatial resolution varies across datasets, with high-resolution products such as Sentinel-1 complementing coarser-scale inputs such as SMAP.

\subsubsection*{Weather Data}
Weather variables were obtained from the Open-Meteo historical archive API \cite{openmeteo}. This source was selected due to its availability of hourly observations without requiring extensive interpolation.

The primary variables used were rain and total precipitation (in millimeters). Hourly values were aggregated into daily totals before merging with station-level observations.

\subsubsection*{Ground Truth Soil Moisture Data}
Ground-truth measurements were collected from in-situ monitoring stations. The temporal modeling experiments initially focused on five stations in Washington State.

For spatial generalization experiments, the dataset was expanded to include additional stations in climatically similar regions, including Idaho and Oregon.



For temporal modeling experiments, five stations in Washington State were used: Darrington-21-NNE and Quinault-4-NE (USCRN), and Spokane-17-SSW (USCRN), SourdoughGulch, and Touchet (SNOTEL). These stations span the state's major climate regimes, from the wet maritime conditions of the western Cascades and Olympic Peninsula to the semi-arid continental interior of eastern Washington.

The response variable for all supervised learning tasks is soil moisture measured at 5 cm depth.

\subsubsection*{External Test Datasets}

Two types of external datasets were considered for evaluating generalization to new locations.
The first consists of unseen public stations outside Washington with similar geographic and climatic characteristics. These datasets are used to evaluate temporal-spatial transfer.
The second consists of ECE sensor deployments within Washington, which provide an additional real-world evaluation setting outside the original training stations.

\textbf{TODO (ECE)}

\subsubsection*{Temporal Coverage}

\begin{table}[H]
\centering
\begin{tabular}{ll}
\toprule
\textbf{Stage} & \textbf{Time Range} \\
\midrule
Initial target window & 2008 -- 2025 \\
Aligned data window & 2011 -- 2025 \\
Final study window (SMAP-constrained) & 2017 -- 2025 \\
\bottomrule
\end{tabular}
\caption{Evolution of the temporal coverage as data sources were integrated.}
\end{table}

The final study period (2017-2025) ensures consistent availability of all features, particularly SMAP-derived variables.
